\chapter{纯化}
\chapterkana{純化}
\epigraphbox{“纯粹而不杂，静一而不变。”}{——《庄子·刻意》}

\kaishou{操作表上有四列。前三列她认得。\par
最后一列她不认得，那一列是退出。\par
她按了前三列。}

\section{处置的通路}

新历的条文顺着上行信道传回月面之后，月之都在第二十七枚边栏写下了「底牌见于人手」六个字，紧接着便定下了处置方案。城内既没有军队，工籍四类杂务全由月兔承担，城中技术又专为无秽之身设计（《章程总目》六卷里连医卷都不设），因此月之都手里没有别的工具，唯一能送到地面去的，依然是历法里那条双向信道。信道在城地之间的运转方式已于第七章第四节交代过，这里只看那道处置令是如何发出的。

发令的规矩载于第四十枚，全文不长。

\begin{archive}[内部参考]{档案 № ФВ-1971-22}
第四十枚。凡处置之令，都由信道下行；一人手书；不署受令者。\\
边栏另有一行，刀迹不同：「令出，不问其行。」
\end{archive}

手书发令之人便是嫦娥。她自公元前二〇五〇年下地起便一直守在这一端，至此时已守了整整六十年。令上之所以不署受令者姓名，也实在不必写：当时地面上依历行事的核心人物唯独后羿一人，条文发下去自然便落到他手上；发令之后不问执行过程，也是同一个道理。

不过执行的速度变化，还是在上行录里留下了痕迹。此前三年，刑戮一类条文从发出到办毕平均耗时十四日；此后一年骤然缩至三日，再过一年竟缩至二日。同一条信道上的农时与祭祀类条文耗时仅少了一两日，唯独刑戮一类提速如此之快。同批残卷中另有一枚记作通例：「凡令至，其人自办，不知其故；事后问之，答曰，历上如此。」那一道处置令发下来时，格式与寻常吉凶宜忌并无二致，日子在前，宜忌在中，名字缀在末尾。

\begin{archive}[内部参考]{档案 № ФВ-1971-22}
第一年的第九十日，宜刑。\\
受令者，伯封。
\end{archive}

伯封是后羿的儿子，处置的字样就这样挂在他的名字下方，顺着历法条文送往了地面。《左传》昭公二十八年记伯封「实有豕心，贪惏无餍，忿颣无期」，终为后羿所灭，父子相搏，同归于尽；《楚辞·天问》则只问不答，留下「何羿之射革，而交吞揆之」一问。\footnote{《左传》昭公二十八年；《楚辞·天问》。本书取伯封为后羿之子一说，理由见本章末。}传世文献给的都是惨烈的结局，至于信道条文在送达那一刻究竟在人的意识里起了何种作用，档案中只有一串编号而全无说明，本书也就不再深挖。

\section{纯狐}

纯狐是后羿之妻、伯封之母，有穷氏的地上土著。她不识字，也不通术数，生平从未登上月球，更未得到过月之民传授的任何技术，手上有的只是三样寻常之物：一部旧历抄本、一处地下遗迹，以及遗迹外壳上的一串刻痕。

《天问》里留有一句「浞娶纯狐，眩妻爰谋」，把她记作了寒浞之妻。\footnote{《楚辞·天问》。这一句把她记作寒浞之妻；本书取「后羿之妻」一说，见本章末。}这六个字之外，先秦文献提及她之处寥寥无几，后世故事则多半渲染她嫉妒狠辣、一心报复私仇。按本书的读法，那其实是一场处置令落回处置者自家的反噬：条文顺着信道下行，后羿依令行事，到头来死的却是她的丈夫与儿子。

照常理，寻仇应当找下令的那一头，可下令的月之都远在月面，不在有穷氏的疆土之上。她所能触及的最远处，便是信道在地面这一头的枢纽，而那一头恰好守着嫦娥。嫦娥在地面常驻了六十年，挂名「地之主任」，在当地算不得秘密，凡有下行条文也必经她手；因此纯狐「找上她」并不难理解，真正难解的是，一个不识字的地上女子最初是如何锁定这一目标的。

这里可以摆出三种猜测，本书不替她强选其一。其一，旧历抄本末尾那一段记号本就是嫦娥一系带下来的，纯狐由记号认出了出处；其二，历法条文虽不署名，六十年的行文习惯却足以泄露出是谁的手笔；其三，她根本不知道对面坐着谁，只是顺着那条通路一路往上追，而往上走原本就只有这一个方向。笔者个人略倾向前两种可能里的某一种，却拿不出任何材料把这三种可能拆开，这一段空白比前几章都要大得多。

\section{旧历抄本}

纯狐手上的那部旧历，是本章唯一可以直接捧在手里读的实物。它并未收入基金会早期的月面编目，而是在一九七二年那次勘察之前由当地村民交出来的，其中夹着一束竹片，长约二十三厘米、宽约一厘米，原本以麻绳编连，出土时绳已朽断。竹片上书写着日序与节气，字迹属于楚系写法，与中原通行的秦篆并非一路。\footnote{这一束现编为 \texttt{ФВ-1972-07}，存三十一片，其中十八片字迹可辨。释文未刊，本书用的是释读稿。}

细看这束竹片，有三处特征值得记上一笔。首先，它确系抄本无疑，竹片上有两三处涂改添补，新笔画直接压在原文之上，显然是抄手一边抄录一边修正自己抄错的字句，与原典内容无关。其次，它的日序与中原通行历法差着几天，大概每过几十天便差出一天上下，逐年累积，这正符合删去三术后的劣化历特征。最后，也是最关键的一点，这一束竹片的末尾有六七片的内容与日序毫无关系，上面一列一列排着一些不成文句的记号，而这几片记号，正与第二章那处地下遗迹外壳上的刻痕属于同一套符号。

将这三点串在一处，纯狐手里的筹码便一清二楚了：一部越走越不准的旧历，外加一段她读不懂的符号。她不需要知道那串符号在文法上作何解释，只要认得出遗迹石壁上刻着一模一样的图形便已足够，这个「一模一样」便是她全部行动的起点。

\section{设备与操作表}

要动用那套记号，还得有执行它的器物。一九七二年测绘过的那处有穷氏地下遗迹（档案 \texttt{ФВ-1972-06}）在第二章已经引过：深约十四米，长约六十米，宽约二十米，四壁规整，底部是一层金属反射面，腔壁上留有四列不可识读的刻痕。第二章仅用它来证明地上设施与天上月之都同源，未曾细究那四列刻痕；本章要解开的，正是这四列刻痕的用途。

四列之中，最右侧一列排列最密，字距约为其余三列的三分之一。按后世对同类遗存设备的整理比对，这四列刻痕其实是一张完整的操作表：自左向右前三列依次为定位、启动与监视，唯有最右一列是退出顺序。设备停止运转之前，操作者必须依照最右一列逐步退出工作范围，人才能安然脱身。换言之，前三列讲的都是「对目标做什么」，唯独最右一列讲的是「操作者怎么出来」，四列之中只有它关乎操作者自身的性命。

至于工作范围如何划定，《规程残卷》第十枚里留有一句简短的通则：「凡器，以执器者为心，周百步，都在其内。」按秦汉尺度，百步折合约一百三十八米，与地下空腔及其周边尺度大体相合。\footnote{规程残卷中这一条不在与地面有关的那四枚里，出自第十枚；本书用它，是因为它讲的是器的通例，不限于月面。}这套机制最致命之处在于它从不区分敌我，凡在工作范围之内的一切物质都一视同仁地加以处理，自然也包括正在按表操作的那个人自己。不过这台设备今日已然无存：一九七二年测绘时底部金属反射面尚在，内部已无可运转的机械结构；八十年代再次复查时，整个空腔均已塌陷。今日所能确认的，只是它当年确曾运转过，而今已成一堆废墟。

\section{划定、启动、退出}

一个不识字的地上女子究竟是如何学会操纵那前三列的，现有材料给不出直接答案，笔者只能顺着三样已知物件拼出一条可能的路径：她幼年在遗迹旁见过刻痕并记住了图形；后来在旧历抄本末尾几片上重逢了完全相同的符号（辨认图形的异同本不需要识字）；最后又摸到了遗迹底下的那台设备。图形、顺序与机器三者凑齐，她便握有了一件别人毫无头绪的工具。

根据事后留下的结果反推，她当时依次做了三步：划定工作范围、启动设备，以及没有退出。前两步尚有遗迹与受创部位可作物证，第三步「没有退出」则完全是从结局倒推出来的。她当年为何没有执行最右那一列退出程序，档案未留一字，大概不外乎三种猜想：要么她根本不识字、读不懂最密的那一列退出指令；要么启动之后已来不及撤步；要么她压根儿不知道前面三列之外还有第四列保命的规程。这三种可能本书同样不替她择一而定，能确定的唯有冰冷的结果：设备启动之后，工作范围内的一切物质都被拆解处理，人体一旦落入此列，便不再是一个活着的整体，她的死也正源于此。

\section{纯化之后}

经此一击，嫦娥身上被处理的那一片材料依然停留在原处，质量未减，位置未变，只是在属性上不再属于她自己。蓬莱人的再生机制依赖一个根本前提，即愈合的对象必须被机体识别为「自身的一部分」；一旦这一前提被不可逆地拆解，那一片组织便既不愈合也不腐烂，彻底脱离了活体的管理。

更麻烦的是，这股纯化作用并未就此止步在最初那一片上。由于设备未曾按正常程序关闭，残余作用的界线一直在缓缓向外推移。起初受影响的只有一小片躯体，其余部分仍属她本人；每当界线之外的健康组织试图再生修补时，新生的组织便以那片已脱离的物质为参照，随即也被划入界线之内。

月之都档案对这桩变故只留了一行字，末尾只注着「未解」二字。若要猜度档案执笔者当年为何写下「未解」，大致有三层顾虑叠在一处：其一，他们认得出这是自家设备的路数，却实在想不通为何会用在这个方向；其二，那台设备两万年间在他们的认知里只用来加工材料，从未有过作用于活体的先例；其三，也是最写不出口的一层：一旦在正式档册里写明真相，便等于亲口承认一个地上土著不仅两次接触到了月之都的核心装备，第二次甚至还将其用在了月之民自己身上。于是档册上便只剩了那一行记录，外加「未解」两个字。至于这三层顾虑之后嫦娥如何熬过一千七百六十四年、她的隔离记录又是怎样一栏一栏写下去的，将由第九章接着交代。

平心而论，这一章有三处全靠推理接续：她当年究竟是如何读懂那张操作表的、她为何没能退出工作范围，以及那一片创面为何始终无法合拢。其中第一处最为要害，偏偏手头一点材料也没有：三列刻痕、旧历末几片的形状与幼年记下的图形，这三样东西究竟如何在一个不识字的人手里凑成一次精确的操作，本书目前还拼不出一条严密的链条，只能算是一个尚待填补的猜测。至于《天问》将纯狐记作寒浞之妻，本书却偏取「后羿之妻」一说，理由全在伯封身上：唯有伯封是后羿之子，那道下行处置令的因果才落得到他头上；若纯狐改嫁寒浞，这层报应的链条便立刻断了。两说各有文献依据，本书选的不过是能与月面档案咬合的那一条；至于是不是选错了文献，这个风险笔者认。
